% =====================================================================
%  Chapter 29 — The Tax Engine
% =====================================================================

\chapter{The Tax Engine}
\label{ch:tax-engine}

The book runs two portfolios in parallel: one taxable, one tax-exempt,
with identical weight paths at every month end. This chapter describes
the engine that produces the difference between them, and the two rules
in it that make the difference behave correctly.

\section{The account model}

An account holds cash and a set of \emph{lots}. A lot is a purchase of a
quantity of one security, with an acquisition date and a cost basis. When
the portfolio is sold, the engine disposes of lots first-in-first-out, and
each disposal produces a gain or loss, a holding period, and a tax.

The account also holds two running balances: the cumulative capital loss
position, which carries forward and offsets capital gains, and the
cumulative net capital loss deducted against ordinary income, which is
capped at \$3,000 a year.

At each year end, the engine settles: it computes the year's realised
gains and losses, applies the capital-loss offset rule, computes the tax
on long-term gains at the preferential rate, on short-term gains at the
ordinary rate, and on qualified dividends at the preferential rate plus
the net investment income tax, deducts the capital loss against ordinary
income up to the cap, carries the rest forward, and pays.

The two accounts run the same weights and the same trades. The tax-exempt
account does everything except pay tax. The difference between their
values at any date is therefore the accumulated cost of the tax code, and
it is a clean measurement of exactly one thing.

\section{Rule one: the holding period}

A lot held for more than one year is long-term; one year or less is
short-term. The engine computes the holding period from the acquisition
date to the disposal date.

The rule interacts with the wash-sale machinery in a way that is easy to
get wrong and that the engine handles explicitly: when a loss is
disallowed because of a replacement purchase, the replacement lot
inherits the original lot's holding period. A wash sale therefore does
not merely defer the loss --- it can convert a later short-term gain into
a long-term one, which is worth real money.

\section{Rule two: the wash sale, in both directions}

Under the wash-sale rule, a loss is disallowed if the investor acquires
the same or a substantially identical security within thirty days before
or after the sale.

Three consequences follow, and the engine implements all three.

\textbf{It operates on purchase.} A purchase of a security within thirty
days after a loss was realised on that security disallows the loss, even
if the purchase is made by a different account. This is why the book's
two accounts cannot be treated as independent: the taxable account's loss
and the exempt account's purchase interact. The engine models the wash
sale across the pair of accounts, which is the conservative treatment.

\textbf{It operates on sale.} A sale within thirty days after a purchase
that was itself a replacement for an earlier loss also disallows the
loss, and the basis and holding period carry over.

\textbf{The basis is adjusted.} The disallowed loss is added to the basis
of the replacement shares. A lot bought for \$100 after a \$10 disallowed
loss has a basis of \$110, and its holding period is the original's.

\begin{takeaway}[Why the rule is modelled in both directions]
A system that only checks the sale side of the wash-sale rule will
systematically overstate after-tax returns, because it will believe that
it can harvest a loss and then immediately buy the same fund back. The
book's engine tracks the thirty-day window on both sides, across both
accounts, and adjusts the basis and the holding period.
\end{takeaway}

\section{The tax schedule}

The book uses the 2025 rules for a single taxpayer with high taxable
income.

\textbf{Long-term capital gains} are taxed at fifteen per cent, the
maximum preferential rate, with the ordinary zero and twenty per cent
brackets applying below that.

\textbf{Short-term capital gains} are taxed as ordinary income.

\textbf{Qualified dividends} are taxed at fifteen per cent, with the
3.8 per cent net investment income tax added for a taxpayer above the
threshold.

\textbf{Ordinary dividends} not qualifying are taxed as ordinary income.

\textbf{Capital losses} offset capital gains without limit, and offset
ordinary income up to \$3,000 a year, with the unused portion carried
forward.

The engine's schedule is a parameter of the system, and the taxable and
exempt accounts are the two ends of a comparison in which the only
difference is the payment of tax.

\section{What the engine deliberately does not do}

Three exclusions, for the same reason the feature layer excludes
fundamental data: each is a place where a wrong assumption would be
invisible in the results.

It does not model state and local income tax, or the deductibility of
investment expenses, or the alternative minimum tax, or the qualified
dividend and capital-gain rates' interaction with the brackets above
fifteen per cent. Each would change the tax bill by some fraction of a
percentage point and each would add a parameter that could not be
checked.

It does not model the wash sale for \emph{substantially identical} securities
in the strict sense --- that is, it treats a different fund tracking the
same index as a different security, which is correct: the Internal
Revenue Service treats a purchase of a substantially identical security as
one that is different in name but the same in substance, and the
substantial-identity standard is a factual question the book cannot
resolve from a price series. The engine is therefore \emph{optimistic}
about wash sales across near-duplicate funds, and says so. The effect is
bounded: the funds that are near-duplicates in this universe are few, and
the system's turnover is low.

It does not model the tax cost of the fund's own internal trading. That
cost is real and it is the reason ETF structures distribute so few capital
gains; the book reports the structural statistics in Chapter 10 rather
than charging an estimate it cannot compute.

\begin{keytakeaways}
\item The engine holds lots with acquisition dates and bases, disposes
first-in-first-out, and settles annually with the capital-loss offset,
the carryforward and the preferential rates.
\item The wash sale is modelled in both directions, across both accounts,
with the basis adjustment and the holding-period tack.
\item The parallel taxable and tax-exempt accounts have identical weight
paths, so the gap between them measures exactly one thing: the cost of
the tax code.
\item The engine omits state tax, the alternative minimum tax and the
substantially-identical-security standard, and the book says so rather
than charging an estimate it cannot check.
\end{keytakeaways}
